
In $\PG(n,q)$, an arc is a set of points, no $n+1$ in a hyperplane.
A cap is set of points, no three collinear.
Here, we restrict our attention to arcs in $\PG(2,q).$  
Arcs in higher dimensional projective spaces 
are equivalent to MDS codes 
and will be treated in Section~\ref{sec:coding}.
Our main examples will be the construction of the Lunelli-Sce hyperoval in $\PG(2,16)$ 
(cf.~\cite{LunelliSce58})
and the Pellegrino cap in $\AG(4,3).$ 
The uniqueness of this cap was proven by Hill~\cite{Hill83}.

\bigskip


A $(k,d)$-arc in a projective plane $\pi$ 
is a set $S$ of $k$ points 
such that any line intersects $S$ in at most $d$ 
points. 
Two arcs $S_1$ and $S_2$ 
are equivalent if there is a projectivity 
$\Phi$ such that $\Phi(A) = B.$ 
The problem of classifying arcs 
is the problem of determining the orbits of the projectivity group on arcs. 
At times, 
we consider the larger group of collineations. 
In that case, the problem of classifying arcs is the 
problem of determining the orbits of the collineation group on arcs. 
Table~\ref{tab:arcs} list the commands available to classify arcs.
\orbitertableofcommands{arc_generator}
Here is an example. 
A hyperoval in a finite plane $\PG(2,2^e)$ is a $(2^e + 2,2)$-arc. 
It is an interesting but difficult problem to classify the hyperovals in finite 
projective planes under the appropriate group action.
For small planes, computer search with Orbiter can help.
The command 

\sourcecode{hyperoval_16_classify}

performs the classification of hyperovals in $\PG(2,16).$ 
There are exactly two hyperovals in this plane. 
Orbiter also finds the stabilizers of these arcs. 
They have orders $16320$ and $144,$ respectively.
The two hyperovals are the regular hyperoval and the Lunelli-Sce hyperoval. 
Here is the relevant output from the Orbiter report 
(in the output, the Lunelli-Sce hyperoval is orbit 0, and the 
regular hyperoval is orbit 1):
\orbiteroutput{GRAPHICS/LATEX/arc_Lunelli_Sce}





\bigskip

In the theory of cubic surfaces, 
we are interested in non-conical arcs of size 6. 
Here, non-conical means that the set of points does not lie on a conic.
Cubic surfaces are associated with non-conical arcs of size 6 
(in a many-to-one relationship when considering isomorphism classes).
The following example demonstrates 
how non-conical 6-arcs can be classified in Orbiter:

\sourcecode{nc_arcs_16}

The number of Eckardt points of the surface can be recovered from properties of the arc.
For this reason, it is interesting to classify arcs so that the associated 
cubic surface has a fixed number of Eckardt points.

